\documentclass[11pt,a4paper]{article}
\usepackage[margin=23mm,headheight=15pt,headsep=8mm,footskip=11mm]{geometry}
\usepackage{amsmath,amsthm}
\usepackage{fontspec}
\usepackage{unicode-math}
\setmainfont{TeX Gyre Pagella}
\setsansfont{TeX Gyre Heros}
\setmonofont{Latin Modern Mono}[Scale=MatchLowercase]
\setmathfont{TeX Gyre Pagella Math}
\usepackage{microtype,booktabs,array,enumitem,xcolor,fancyhdr,needspace}
\usepackage{hyperref}
\definecolor{inkblue}{HTML}{243F56}
\definecolor{rulegray}{HTML}{AEB7BE}
\hypersetup{colorlinks=true,linkcolor=inkblue,citecolor=inkblue,urlcolor=inkblue,
  pdftitle={Rank-three finite T-data: a computer-assisted classification},
  pdfsubject={Sixteen classes, explicit matrices, proofs and exact certificates},
  pdfauthor={}}
\pagestyle{fancy}
\fancyhf{}
\fancyhead[L]{\small\sffamily Rank-three finite T-data}
\fancyhead[R]{\small\sffamily Classification and certificates}
\fancyfoot[C]{\small\thepage}
\renewcommand{\headrulewidth}{0.3pt}
\setlength{\parindent}{1.2em}
\setlength{\parskip}{3pt plus 1pt}
\setlength{\emergencystretch}{1em}
\setlist{itemsep=3pt,topsep=5pt}
\raggedbottom
\newtheorem{theorem}{Theorem}[section]
\newtheorem{lemma}[theorem]{Lemma}
\newtheorem{corollary}[theorem]{Corollary}
\theoremstyle{definition}
\newtheorem{definition}[theorem]{Definition}
\newtheorem{computation}[theorem]{Exact computation}
\DeclareMathOperator{\diag}{diag}
\DeclareMathOperator{\Trop}{Trop}
\newcommand{\N}{\mathbb N}
\newcommand{\Q}{\mathbb Q}
\newcommand{\Z}{\mathbb Z}
\newcommand{\R}{\mathbb R}
\newcommand{\code}[1]{\texttt{\detokenize{#1}}}

\title{\vspace{-12mm}\textbf{Rank-three finite T-data}\\[4pt]
  \large A computer-assisted classification}
\author{}
\date{5 September 2026}

\begin{document}
\thispagestyle{plain}
\begin{center}
{\LARGE\bfseries Rank-three finite T-data\par}
\vspace{4pt}
{\large A computer-assisted classification\par}
\vspace{9pt}
{\small 5 September 2026}
\end{center}
\begin{abstract}
In the setting of identity symmetrizer and diagonal leading matrix, the
classification consists of sixteen indecomposable families. Every finite-type
datum is obtained from one of sixteen explicit representatives by an admissible
rational rescaling of time and shifts of the three species, together with
relabelling and sign exchange. The completeness argument combines a
delay-independent bound on coefficient sums with exact integer feasibility
checks. Two-sided reddening certificates establish periodicity of all
representatives. In this setting, simultaneous positivity is both necessary
and sufficient for finite type in rank three.
\end{abstract}

\section{Statement of the classification}\label{sec:statement}

The scope is the one used in Mizuno's rank-two classification
\cite{MizunoRankTwo}: the symmetrizer is $D=I_3$, and the leading matrix is
diagonal. The general definition of T-data in \cite{MizunoTData} also permits
nontrivial symmetrizers and non-diagonal leading matrices.

\begin{definition}\label{def:datum}
Let $r_1,r_2,r_3$ be positive integers, and let $N_+(z),N_-(z)$ be
$3\times3$ matrices over $\N[z]$. Set
\begin{equation}\label{eq:leading}
 N_0(z)=\diag(1+z^{r_1},1+z^{r_2},1+z^{r_3}),\qquad
 A_\pm(z)=N_0(z)-N_\pm(z).
\end{equation}
In entry $(i,j)$, all nonzero coefficients of $N_\pm$ occur at exponents
$0<p<r_i$, and the supports of the two signs are disjoint. The pair is a
T-datum in the present setting if
\begin{equation}\label{eq:symplectic}
 A_+(z)A_-(z^{-1})^{\mathsf T}
 =A_-(z)A_+(z^{-1})^{\mathsf T}.
\end{equation}
It is \emph{indecomposable} if no simultaneous permutation of the indices
makes both matrices block diagonal with two nonempty blocks.
\end{definition}

Write $n_{ij;p}=[z^p](N_+-N_-)_{ij}$ and $[a]_+=\max(a,0)$.
The associated Y-system is
\begin{equation}\label{eq:ysystem}
 Y_i(u)Y_i(u-r_i)=
 \prod_{j=1}^{3}\prod_{p=1}^{r_i-1}
 Y_j(u-p)^{[n_{ij;p}]_+}(1+Y_j(u-p))^{-n_{ij;p}}.
\end{equation}
Finite type means periodicity of this system in the universal semifield,
equivalently the condition in \cite[Definition 1.4]{MizunoRankTwo}.

\begin{theorem}[Computer-assisted classification]\label{thm:classification}
An indecomposable T-datum in Definition~\ref{def:datum} is of finite type if
and only if, after relabelling the indices and possibly exchanging the signs,
it has the form
\begin{equation}\label{eq:normal-form}
 A_\pm(z)=E(z)B_\pm^{(k)}(z^\lambda)E(z)^{-1},\qquad
 E(z)=\diag(z^{s_1},z^{s_2},1),
\end{equation}
for exactly one of the sixteen representatives in
Section~\ref{sec:catalogue}. Here $\lambda\in\Q_{>0}$ and $s_1,s_2\in\Q$;
the parameters must make every exponent integral and satisfy the support
conditions in Definition~\ref{def:datum}.
\end{theorem}

The expression in \eqref{eq:normal-form} is first interpreted in the group
algebra of rational powers of $z$. If a representative has delay $r_i^{(k)}$
and a term $z^p$ in entry $(i,j)$, the transformed delay and exponent are
\begin{equation}\label{eq:parameters}
 r_i=\lambda r_i^{(k)},\qquad p'=\lambda p+s_i-s_j,\qquad s_3=0.
\end{equation}
Thus the statement gives all polynomial lifts, including arbitrarily large
delays. Admissible rescalings and shifts correspond to changes of slices;
the sixteen slice classes are distinct, as discussed in
Section~\ref{sec:slices}.

\begin{corollary}\label{cor:positivity}
An indecomposable rank-three T-datum in this setting is of finite type if
and only if there is a row vector $v\in\R^3$ such that
\begin{equation}\label{eq:positivity}
 v>0,\qquad vA_+(1)>0,\qquad vA_-(1)>0.
\end{equation}
All inequalities are componentwise.
\end{corollary}

Necessity is the simultaneous positivity theorem
\cite[Theorem 5.5]{MizunoTData}, recalled in
\cite[Lemma 3.3]{MizunoRankTwo}. Sufficiency follows from the exhaustive
reduction and the periodicity certificates below. Decomposable rank-three
data are direct sums with rank partitions $2+1$ or $1+1+1$, using the
rank-two and rank-one classifications.

\paragraph{Proof status.}
The proof uses exact integer and rational arithmetic and the UNSAT results
of Z3. All 180 unbounded solver queries are saved and replayed. The
certificate computation is reproducible; it has not been formalized in a
proof assistant. The source data and runtime details are given in
Appendix~\ref{app:reproduce}.

\section{Periods and the structure of the proof}\label{sec:overview}

Table~\ref{tab:periods} records the reddening lengths $h_+,h_-$ in the
displayed time coordinates, the labelled seed period $\Omega$, and the number of
vertices in one slice. The period is the first return of the principal
C-matrix to the identity in the exact calculation.

\begin{table}[!ht]
\centering\small\renewcommand{\arraystretch}{1.02}
\begin{tabular}{@{}crrrrr@{}}
\toprule
Class & Delays $(r_1,r_2,r_3)$ & $h_+$ & $h_-$ & Period & Slice size\\
\midrule
1 & $(2,2,2)$ & 2 & 4 & 12 & 3\\
2 & $(2,2,2)$ & 2 & 7 & 9 & 6\\
3 & $(3,3,2)$ & 3 & 8 & 11 & 8\\
4 & $(6,2,6)$ & 6 & 14 & 20 & 7\\
5 & $(2,2,8)$ & 12 & 12 & 24 & 6\\
6 & $(2,2,12)$ & 12 & 24 & 36 & 8\\
7 & $(5,2,2)$ & 5 & 7 & 12 & 9\\
8 & $(8,2,6)$ & 8 & 14 & 22 & 8\\
9 & $(5,2,3)$ & 5 & 8 & 13 & 10\\
10 & $(8,2,6)$ & 8 & 16 & 24 & 8\\
11 & $(2,2,8)$ & 10 & 8 & 18 & 6\\
12 & $(2,2,12)$ & 14 & 18 & 32 & 8\\
13 & $(3,2,2)$ & 3 & 7 & 10 & 7\\
14 & $(2,12,10)$ & 22 & 18 & 40 & 12\\
15 & $(2,2,2)$ & 4 & 3 & 7 & 6\\
16 & $(2,2,2)$ & 4 & 3 & 14 & 6\\
\bottomrule
\end{tabular}

\caption{Periods of the sixteen representatives.}\label{tab:periods}
\end{table}

The proof has three parts. First, positivity bounds each off-diagonal
coefficient sum by $7$, reducing the problem to $180$ constant-matrix pairs
up to relabelling and sign exchange. Second, exact linear integer constraints
exclude $164$ pairs and classify every lift of the remaining $16$ pairs.
Third, positive and negative reddening sequences prove finite type for
each representative, and rescaling and shifting transfer periodicity to
its entire family.

\section{The reduction to finitely many constant pairs}\label{sec:constants}

Throughout this section write $P=N_+(1)$ and $M=N_-(1)$. If the datum is of
finite type, \eqref{eq:positivity} gives
\begin{equation}\label{eq:pm-positive}
 v(2I-P)>0,\qquad v(2I-M)>0.
\end{equation}
The diagonal entries of $2I-P$ and $2I-M$ are positive integers at most $2$.
Hence $P_{ii},M_{ii}\in\{0,1\}$. Both Z-matrices are nonsingular M-matrices,
so all their principal minors are positive; this also follows by diagonal
scaling using $v$.

\begin{lemma}[Strong connectivity]\label{lem:strong}
For every indecomposable datum in Definition~\ref{def:datum}, the directed
graph with an arrow $i\to j$ when $i\ne j$ and $P_{ij}+M_{ij}>0$ is strongly
connected. This assertion does not require finite type.
\end{lemma}
\begin{proof}
Suppose a cut makes both $A_+$ and $A_-$ block lower triangular. Their
constant terms are identity matrices. Thus $K=A_+^{-1}A_-$ exists over
$\Q(z)$ and is regular at zero. The symplectic identity implies
\[
 K(z)=K(z^{-1})^{\mathsf T}.
\]
The block lower triangular matrix $K$ is therefore block diagonal.

If a crossing block were nonzero, let $h>0$ be the smallest degree of
an entry crossing the cut. The identity
\[
 K=I+A_+^{-1}(N_+-N_-)
\]
shows that the degree-$h$ crossing coefficient of $K$ equals the degree-$h$
crossing coefficient of $N_+-N_-$. A term involving a nonconstant factor
of $A_+^{-1}$ has larger degree, since a crossing factor already has degree
at least $h$. Disjoint support makes this coefficient nonzero, a
contradiction. A reducible support graph consequently has no edges between
its strongly connected components, and the datum is decomposable.
\end{proof}

\begin{lemma}[Coefficient bound]\label{lem:bound}
For an indecomposable finite-type datum, set
$w_{ij}=\max(P_{ij},M_{ij})$ for $i\ne j$. Then
\begin{align}
 w_{ij}w_{ji}&<4,\label{eq:two-cycle}\\
 w_{12}w_{23}w_{31}&<8,\qquad w_{13}w_{32}w_{21}<8.\label{eq:three-cycle}
\end{align}
In particular, $0\le P_{ij},M_{ij}\le7$ for $i\ne j$.
\end{lemma}
\begin{proof}
For every nonzero $w_{ij}$, \eqref{eq:pm-positive} gives
$w_{ij}v_i<2v_j$. Each arrow in a strongly connected three-vertex graph
belongs to a simple directed cycle of length two or three. Multiplication
along that cycle proves the displayed inequalities. The other factors in
a nonzero cycle product are positive integers, giving the bound $7$.
\end{proof}

The program \code{enumerate_constants.py} exhausts the resulting finite
set. It enumerates the six maxima $w_{ij}$, every ordered split
$(P_{ij},M_{ij})$ with that maximum, and all $64$ choices of the diagonal
entries. It imposes
\[
 (2I-P)(2I-M)^{\mathsf T}=(2I-M)(2I-P)^{\mathsf T},
\]
positive principal minors and exact simultaneous positivity. The counts are
shown in Table~\ref{tab:counts}. The bounds restrict coefficient sums;
no upper bound is imposed on the delays.

\begin{table}[ht]
\centering\small\renewcommand{\arraystretch}{1.15}
\begin{tabular}{@{}p{0.79\linewidth}r@{}}
\toprule
Stage & Count\\\midrule
Strongly connected maxima satisfying the cycle inequalities & 1,134\\
Ordered off-diagonal splits & 1,474,290\\
Labelled pairs after symplectic and principal-minor tests & 1,835\\
Orbits under relabelling and sign exchange & 180\\
Orbits satisfying simultaneous positivity & 180\\
\bottomrule
\end{tabular}
\caption{The exhaustive constant-matrix reduction.}\label{tab:counts}
\end{table}

For the positivity test, the program considers the cone given by the nine
weak inequalities $v\ge0$, $v(2I-P)\ge0$ and $v(2I-M)\ge0$. It is pointed
because it lies in the nonnegative orthant. In dimension three its extreme
rays arise from cross products of two constraint normals. Summing the
admissible rays and testing the sum against the strict inequalities is an
exact feasibility test. An explicit positive integer vector is saved for
every retained pair.

\section{All polynomial lifts, with unbounded delays}\label{sec:lifts}

For each of the $180$ constant pairs, introduce an integer variable for
every delay and for each unit of coefficient mass in every entry.
Units of the same sign may have equal exponents, so coefficients greater
than one are included. Exponents of opposite signs in the same entry must
differ. Ordering exponents within each sign and entry removes duplicate
labellings.

For a matrix $F$ over $\Q(z)$, write $F^*(z)=F(z^{-1})^{\mathsf T}$.
Each entry of $A_+A_-^*-A_-A_+^*$ expands as a difference of two finite
multisets of Laurent monomials. Each exponent is an integral linear form
in the delay and exponent variables. The polynomial identity is exactly
equality of these exponent multisets.

The solver equates multiplicities at every exponent appearing on the
positive side. Equality of the total multiplicities follows from the
constant-matrix identity. This encoding is exact even when monomials
coalesce. The resulting formulas belong to linear integer arithmetic with
Boolean combinations of equalities, with no delay or exponent cutoff.

\begin{computation}\label{comp:lifts}
Of the $180$ constant pairs, $164$ have no admissible polynomial lift.
The remaining $16$ are the pairs in Section~\ref{sec:catalogue}. For each
survivor, all admissible lifts lie in one three-dimensional rational linear
subspace generated by its representative's exponent vector and the two
relative species shifts in \eqref{eq:normal-form}.
\end{computation}

For each survivor, \code{classify_lift_spaces.py} pairs equal exponents
in a satisfying solution. This matching gives a homogeneous rational
linear subspace on which the entire Laurent identity holds. The program
then asks whether an admissible solution exists outside that subspace.
All sixteen answers are UNSAT, and each subspace has dimension exactly
three. Its independent generators are the representative's exponent
vector and the two shifts $s_1,s_2$.

The verifier separately reduces every exponent modulo the row-reduced
linear equations and checks equality of the resulting monomial multisets.
Thus the symplectic identities hold symbolically throughout the families.
It also replays all $164$ exclusion queries and all $16$ queries for a
solution outside the family. The complete SMT-LIB2 formulas and their
SHA-256 hashes are saved. Every replayed answer is UNSAT.

This proves necessity in Theorem~\ref{thm:classification}: finite type
implies one of $180$ constant pairs, and every admissible symplectic lift
of such a pair belongs to one of the sixteen families.

\section{Periodicity of the representatives and their families}\label{sec:periodicity}

For a delay tuple $\symbf{r}$, the reconstructed exchange matrix has
vertex set
\[
 R=\{(i,p):1\le i\le3,\ 0\le p<r_i\}.
\]
It is constructed from \cite[equation (2.2)]{MizunoRankTwo}. A time step
mutates the three vertices $(i,0)$ and then relabels
$(i,p)\mapsto(i,p-1\bmod r_i)$. The implementation checks skew-symmetry,
commutativity of these mutations, return of the exchange matrix after
every step, and inverse-step consistency.

\begin{computation}\label{comp:reddening}
Starting from the identity principal C-matrix, every representative
reaches a negative permutation matrix in each time direction, at the
lengths $h_+$ and $h_-$ in Table~\ref{tab:periods}. Exact integer mutation
also gives the identity returns recorded in that table.
\end{computation}

The two-sided reddening criterion
\cite[Proposition 2.7]{MizunoRankTwo} proves finite type for each
representative in every semifield. The five examples in the earlier
rank-two table and the sixth rank-two tadpole case reproduce the
reddening lengths of the rank-two classification as controls.

\begin{lemma}\label{lem:invariance}
Finite type is invariant under the admissible transformations
in \eqref{eq:normal-form}.
\end{lemma}
\begin{proof}
For integral species shifts, the transformation is the bijection of
solutions $Y'_i(u)=Y_i(u-s_i)$. Replacing $z$ by $z^m$, for a positive
integer $m$, splits the system into $m$ independent residue classes of
time. It therefore preserves finite type in both directions.

For rational $\lambda,s_1,s_2$, choose a positive integer $L$ clearing
their denominators. The transformed datum at $z^L$ is an integral
species shift of the representative at $z^{L\lambda}$. Apply the integral
argument to both sides.
\end{proof}

The exact reddening certificates and Lemma~\ref{lem:invariance} prove
sufficiency in Theorem~\ref{thm:classification} for every admissible
parameter value. Together with the constant-matrix reduction they also
prove Corollary~\ref{cor:positivity}.

\section{Change of slices and distinctness}\label{sec:slices}

The sixteen constant pairs are distinct under relabelling and sign
exchange. Specialization at $z=1$ is invariant under
\eqref{eq:normal-form}, so its sixteen families do not overlap.

For the slice equivalence of \cite[Section 3.1]{MizunoRankTwo}, integral
time dilations repeat the same slice sequence. Species shifts change
the placement of mutation events and the initial cut. They are connected
by a straight path of admissible rational shifts, since the support
inequalities define a convex domain. The event order changes only at
coincident event times. Clearing denominators makes such a coincidence
a simultaneous mutation in a valid T-datum, so the relevant exchange
entry is zero and the mutations commute. Moving the initial cut rotates
the cyclic sequence. Thus admissible rational rescalings and shifts
give changes of slices.

Distinctness is checked directly on the cyclic mutation sequences.
For each representative, \code{slice_equivalence.py} follows one
connected component through the cyclic permutation of components.
The returning loop has three mutation events and a terminal vertex
permutation. Each event roots one orbit of that permutation, and these
three orbits cover the vertices. Ordering each orbit from its root gives
a canonical vertex labelling. The three orbit lengths and the exchange
matrix in this labelling encode the loop completely.

The program exhausts cyclic rotations and adjacent commuting mutations
until closure. For a loop encoded by $(B,(a,b,c),\pi)$, rotation gives
\[
 (\mu_aB,(b,c,\pi^{-1}a),\pi).
\]
Adjacent events can be exchanged when their exchange entry at that
point vanishes. Sign exchange is $B\mapsto-B$, with the word and
permutation unchanged. The sixteen resulting sets of canonical
encodings are disjoint. The full minimal encoding, as well as its hash,
is saved for every class.

\clearpage
\section{The sixteen representatives}\label{sec:catalogue}
For each class, the constant matrices displayed below are
$A_\pm(1)=B_\pm^{(k)}(1)=2I_3-N_\pm^{(k)}(1)$.
They are unchanged by the admissible rescalings and species shifts in
\eqref{eq:normal-form}; relabelling permutes their rows and columns,
and sign exchange interchanges the pair.
\subsection*{Classes 1--4}
\noindent In each class, $B_\pm^{(k)}(z)=\operatorname{diag}(1+z^{r_1},1+z^{r_2},1+z^{r_3})-N_\pm^{(k)}(z)$.\par\medskip
\noindent\begin{minipage}{\linewidth}
\noindent\textbf{Class 1}\hfill $\symbf{r}=(2,2,2)$
\par\smallskip\noindent\textcolor{rulegray}{\rule{\linewidth}{0.3pt}}
\begin{minipage}[c]{0.495\linewidth}\centering\fontsize{10.5}{13}\selectfont
\setlength{\abovedisplayskip}{3pt}\setlength{\belowdisplayskip}{3pt}
\setlength{\abovedisplayshortskip}{3pt}\setlength{\belowdisplayshortskip}{3pt}
\setlength{\arraycolsep}{3.5pt}\renewcommand{\arraystretch}{1.18}
\[N_{+}^{(1)}(z)=\begin{pmatrix}0 & 0 & 0\\0 & 0 & 0\\0 & 0 & 0\end{pmatrix}\]
\end{minipage}%
\begin{minipage}[c]{0.495\linewidth}\centering\fontsize{10.5}{13}\selectfont
\setlength{\abovedisplayskip}{3pt}\setlength{\belowdisplayskip}{3pt}
\setlength{\abovedisplayshortskip}{3pt}\setlength{\belowdisplayshortskip}{3pt}
\setlength{\arraycolsep}{3.5pt}\renewcommand{\arraystretch}{1.18}
\[N_{-}^{(1)}(z)=\begin{pmatrix}0 & 0 & z\\0 & 0 & z\\z & z & 0\end{pmatrix}\]
\end{minipage}%
\par\noindent
\begin{minipage}[c]{0.495\linewidth}\centering\fontsize{10.5}{13}\selectfont
\setlength{\abovedisplayskip}{3pt}\setlength{\belowdisplayskip}{3pt}
\setlength{\abovedisplayshortskip}{3pt}\setlength{\belowdisplayshortskip}{3pt}
\setlength{\arraycolsep}{3.5pt}\renewcommand{\arraystretch}{1.18}
\[A_{+}(1)=\begin{pmatrix}2 & 0 & 0\\0 & 2 & 0\\0 & 0 & 2\end{pmatrix}\]
\end{minipage}%
\begin{minipage}[c]{0.495\linewidth}\centering\fontsize{10.5}{13}\selectfont
\setlength{\abovedisplayskip}{3pt}\setlength{\belowdisplayskip}{3pt}
\setlength{\abovedisplayshortskip}{3pt}\setlength{\belowdisplayshortskip}{3pt}
\setlength{\arraycolsep}{3.5pt}\renewcommand{\arraystretch}{1.18}
\[A_{-}(1)=\begin{pmatrix}2 & 0 & -1\\0 & 2 & -1\\-1 & -1 & 2\end{pmatrix}\]
\end{minipage}%
\par\smallskip\noindent\small $h_+=2,\quad h_-=4,\quad \Omega=12.$\quad Vertices per slice: 3.
\end{minipage}\par\vfill
\noindent\begin{minipage}{\linewidth}
\noindent\textbf{Class 2}\hfill $\symbf{r}=(2,2,2)$
\par\smallskip\noindent\textcolor{rulegray}{\rule{\linewidth}{0.3pt}}
\begin{minipage}[c]{0.495\linewidth}\centering\fontsize{10.5}{13}\selectfont
\setlength{\abovedisplayskip}{3pt}\setlength{\belowdisplayskip}{3pt}
\setlength{\abovedisplayshortskip}{3pt}\setlength{\belowdisplayshortskip}{3pt}
\setlength{\arraycolsep}{3.5pt}\renewcommand{\arraystretch}{1.18}
\[N_{+}^{(2)}(z)=\begin{pmatrix}0 & 0 & 0\\0 & 0 & 0\\0 & 0 & 0\end{pmatrix}\]
\end{minipage}%
\begin{minipage}[c]{0.495\linewidth}\centering\fontsize{10.5}{13}\selectfont
\setlength{\abovedisplayskip}{3pt}\setlength{\belowdisplayskip}{3pt}
\setlength{\abovedisplayshortskip}{3pt}\setlength{\belowdisplayshortskip}{3pt}
\setlength{\arraycolsep}{3.5pt}\renewcommand{\arraystretch}{1.18}
\[N_{-}^{(2)}(z)=\begin{pmatrix}0 & 0 & z\\0 & z & z\\z & z & 0\end{pmatrix}\]
\end{minipage}%
\par\noindent
\begin{minipage}[c]{0.495\linewidth}\centering\fontsize{10.5}{13}\selectfont
\setlength{\abovedisplayskip}{3pt}\setlength{\belowdisplayskip}{3pt}
\setlength{\abovedisplayshortskip}{3pt}\setlength{\belowdisplayshortskip}{3pt}
\setlength{\arraycolsep}{3.5pt}\renewcommand{\arraystretch}{1.18}
\[A_{+}(1)=\begin{pmatrix}2 & 0 & 0\\0 & 2 & 0\\0 & 0 & 2\end{pmatrix}\]
\end{minipage}%
\begin{minipage}[c]{0.495\linewidth}\centering\fontsize{10.5}{13}\selectfont
\setlength{\abovedisplayskip}{3pt}\setlength{\belowdisplayskip}{3pt}
\setlength{\abovedisplayshortskip}{3pt}\setlength{\belowdisplayshortskip}{3pt}
\setlength{\arraycolsep}{3.5pt}\renewcommand{\arraystretch}{1.18}
\[A_{-}(1)=\begin{pmatrix}2 & 0 & -1\\0 & 1 & -1\\-1 & -1 & 2\end{pmatrix}\]
\end{minipage}%
\par\smallskip\noindent\small $h_+=2,\quad h_-=7,\quad \Omega=9.$\quad Vertices per slice: 6.
\end{minipage}\par\vfill
\noindent\begin{minipage}{\linewidth}
\noindent\textbf{Class 3}\hfill $\symbf{r}=(3,3,2)$
\par\smallskip\noindent\textcolor{rulegray}{\rule{\linewidth}{0.3pt}}
\begin{minipage}[c]{0.495\linewidth}\centering\fontsize{10.5}{13}\selectfont
\setlength{\abovedisplayskip}{3pt}\setlength{\belowdisplayskip}{3pt}
\setlength{\abovedisplayshortskip}{3pt}\setlength{\belowdisplayshortskip}{3pt}
\setlength{\arraycolsep}{3.5pt}\renewcommand{\arraystretch}{1.18}
\[N_{+}^{(3)}(z)=\begin{pmatrix}0 & 0 & 0\\0 & 0 & 0\\0 & 0 & z\end{pmatrix}\]
\end{minipage}%
\begin{minipage}[c]{0.495\linewidth}\centering\fontsize{10.5}{13}\selectfont
\setlength{\abovedisplayskip}{3pt}\setlength{\belowdisplayskip}{3pt}
\setlength{\abovedisplayshortskip}{3pt}\setlength{\belowdisplayshortskip}{3pt}
\setlength{\arraycolsep}{3.5pt}\renewcommand{\arraystretch}{1.18}
\[N_{-}^{(3)}(z)=\begin{pmatrix}0 & z & 0\\z^{2} & 0 & z^{2} + z\\0 & z & 0\end{pmatrix}\]
\end{minipage}%
\par\noindent
\begin{minipage}[c]{0.495\linewidth}\centering\fontsize{10.5}{13}\selectfont
\setlength{\abovedisplayskip}{3pt}\setlength{\belowdisplayskip}{3pt}
\setlength{\abovedisplayshortskip}{3pt}\setlength{\belowdisplayshortskip}{3pt}
\setlength{\arraycolsep}{3.5pt}\renewcommand{\arraystretch}{1.18}
\[A_{+}(1)=\begin{pmatrix}2 & 0 & 0\\0 & 2 & 0\\0 & 0 & 1\end{pmatrix}\]
\end{minipage}%
\begin{minipage}[c]{0.495\linewidth}\centering\fontsize{10.5}{13}\selectfont
\setlength{\abovedisplayskip}{3pt}\setlength{\belowdisplayskip}{3pt}
\setlength{\abovedisplayshortskip}{3pt}\setlength{\belowdisplayshortskip}{3pt}
\setlength{\arraycolsep}{3.5pt}\renewcommand{\arraystretch}{1.18}
\[A_{-}(1)=\begin{pmatrix}2 & -1 & 0\\-1 & 2 & -2\\0 & -1 & 2\end{pmatrix}\]
\end{minipage}%
\par\smallskip\noindent\small $h_+=3,\quad h_-=8,\quad \Omega=11.$\quad Vertices per slice: 8.
\end{minipage}\par\vfill
\noindent\begin{minipage}{\linewidth}
\noindent\textbf{Class 4}\hfill $\symbf{r}=(6,2,6)$
\par\smallskip\noindent\textcolor{rulegray}{\rule{\linewidth}{0.3pt}}
\begin{minipage}[c]{0.495\linewidth}\centering\fontsize{10.5}{13}\selectfont
\setlength{\abovedisplayskip}{3pt}\setlength{\belowdisplayskip}{3pt}
\setlength{\abovedisplayshortskip}{3pt}\setlength{\belowdisplayshortskip}{3pt}
\setlength{\arraycolsep}{3.5pt}\renewcommand{\arraystretch}{1.18}
\[N_{+}^{(4)}(z)=\begin{pmatrix}0 & 0 & 0\\0 & 0 & 0\\0 & z^{3} & 0\end{pmatrix}\]
\end{minipage}%
\begin{minipage}[c]{0.495\linewidth}\centering\fontsize{10.5}{13}\selectfont
\setlength{\abovedisplayskip}{3pt}\setlength{\belowdisplayskip}{3pt}
\setlength{\abovedisplayshortskip}{3pt}\setlength{\belowdisplayshortskip}{3pt}
\setlength{\arraycolsep}{3.5pt}\renewcommand{\arraystretch}{1.18}
\[N_{-}^{(4)}(z)=\begin{pmatrix}0 & 0 & z\\0 & 0 & z\\z^{5} & z^{5} + z & 0\end{pmatrix}\]
\end{minipage}%
\par\noindent
\begin{minipage}[c]{0.495\linewidth}\centering\fontsize{10.5}{13}\selectfont
\setlength{\abovedisplayskip}{3pt}\setlength{\belowdisplayskip}{3pt}
\setlength{\abovedisplayshortskip}{3pt}\setlength{\belowdisplayshortskip}{3pt}
\setlength{\arraycolsep}{3.5pt}\renewcommand{\arraystretch}{1.18}
\[A_{+}(1)=\begin{pmatrix}2 & 0 & 0\\0 & 2 & 0\\0 & -1 & 2\end{pmatrix}\]
\end{minipage}%
\begin{minipage}[c]{0.495\linewidth}\centering\fontsize{10.5}{13}\selectfont
\setlength{\abovedisplayskip}{3pt}\setlength{\belowdisplayskip}{3pt}
\setlength{\abovedisplayshortskip}{3pt}\setlength{\belowdisplayshortskip}{3pt}
\setlength{\arraycolsep}{3.5pt}\renewcommand{\arraystretch}{1.18}
\[A_{-}(1)=\begin{pmatrix}2 & 0 & -1\\0 & 2 & -1\\-1 & -2 & 2\end{pmatrix}\]
\end{minipage}%
\par\smallskip\noindent\small $h_+=6,\quad h_-=14,\quad \Omega=20.$\quad Vertices per slice: 7.
\end{minipage}\par\vfill
\clearpage
\subsection*{Classes 5--8}
\noindent In each class, $B_\pm^{(k)}(z)=\operatorname{diag}(1+z^{r_1},1+z^{r_2},1+z^{r_3})-N_\pm^{(k)}(z)$.\par\medskip
\noindent\begin{minipage}{\linewidth}
\noindent\textbf{Class 5}\hfill $\symbf{r}=(2,2,8)$
\par\smallskip\noindent\textcolor{rulegray}{\rule{\linewidth}{0.3pt}}
\begin{minipage}[c]{0.495\linewidth}\centering\fontsize{10.5}{13}\selectfont
\setlength{\abovedisplayskip}{3pt}\setlength{\belowdisplayskip}{3pt}
\setlength{\abovedisplayshortskip}{3pt}\setlength{\belowdisplayshortskip}{3pt}
\setlength{\arraycolsep}{3.5pt}\renewcommand{\arraystretch}{1.18}
\[N_{+}^{(5)}(z)=\begin{pmatrix}0 & 0 & 0\\0 & 0 & 0\\z^{4} & z^{5} + z^{3} & z^{4}\end{pmatrix}\]
\end{minipage}%
\begin{minipage}[c]{0.495\linewidth}\centering\fontsize{10.5}{13}\selectfont
\setlength{\abovedisplayskip}{3pt}\setlength{\belowdisplayskip}{3pt}
\setlength{\abovedisplayshortskip}{3pt}\setlength{\belowdisplayshortskip}{3pt}
\setlength{\arraycolsep}{3.5pt}\renewcommand{\arraystretch}{1.18}
\[N_{-}^{(5)}(z)=\begin{pmatrix}0 & z & 0\\z & 0 & z\\0 & z^{7} + z & 0\end{pmatrix}\]
\end{minipage}%
\par\noindent
\begin{minipage}[c]{0.495\linewidth}\centering\fontsize{10.5}{13}\selectfont
\setlength{\abovedisplayskip}{3pt}\setlength{\belowdisplayskip}{3pt}
\setlength{\abovedisplayshortskip}{3pt}\setlength{\belowdisplayshortskip}{3pt}
\setlength{\arraycolsep}{3.5pt}\renewcommand{\arraystretch}{1.18}
\[A_{+}(1)=\begin{pmatrix}2 & 0 & 0\\0 & 2 & 0\\-1 & -2 & 1\end{pmatrix}\]
\end{minipage}%
\begin{minipage}[c]{0.495\linewidth}\centering\fontsize{10.5}{13}\selectfont
\setlength{\abovedisplayskip}{3pt}\setlength{\belowdisplayskip}{3pt}
\setlength{\abovedisplayshortskip}{3pt}\setlength{\belowdisplayshortskip}{3pt}
\setlength{\arraycolsep}{3.5pt}\renewcommand{\arraystretch}{1.18}
\[A_{-}(1)=\begin{pmatrix}2 & -1 & 0\\-1 & 2 & -1\\0 & -2 & 2\end{pmatrix}\]
\end{minipage}%
\par\smallskip\noindent\small $h_+=12,\quad h_-=12,\quad \Omega=24.$\quad Vertices per slice: 6.
\end{minipage}\par\vfill
\noindent\begin{minipage}{\linewidth}
\noindent\textbf{Class 6}\hfill $\symbf{r}=(2,2,12)$
\par\smallskip\noindent\textcolor{rulegray}{\rule{\linewidth}{0.3pt}}
\begin{minipage}[c]{0.495\linewidth}\centering\fontsize{10.5}{13}\selectfont
\setlength{\abovedisplayskip}{3pt}\setlength{\belowdisplayskip}{3pt}
\setlength{\abovedisplayshortskip}{3pt}\setlength{\belowdisplayshortskip}{3pt}
\setlength{\arraycolsep}{3.5pt}\renewcommand{\arraystretch}{1.18}
\[N_{+}^{(6)}(z)=\begin{pmatrix}0 & 0 & 0\\0 & 0 & 0\\z^{8} + z^{4} & z^{9} + z^{3} & 0\end{pmatrix}\]
\end{minipage}%
\begin{minipage}[c]{0.495\linewidth}\centering\fontsize{10.5}{13}\selectfont
\setlength{\abovedisplayskip}{3pt}\setlength{\belowdisplayskip}{3pt}
\setlength{\abovedisplayshortskip}{3pt}\setlength{\belowdisplayshortskip}{3pt}
\setlength{\arraycolsep}{3.5pt}\renewcommand{\arraystretch}{1.18}
\[N_{-}^{(6)}(z)=\begin{pmatrix}0 & z & 0\\z & 0 & z\\z^{6} & z^{11} + z & 0\end{pmatrix}\]
\end{minipage}%
\par\noindent
\begin{minipage}[c]{0.495\linewidth}\centering\fontsize{10.5}{13}\selectfont
\setlength{\abovedisplayskip}{3pt}\setlength{\belowdisplayskip}{3pt}
\setlength{\abovedisplayshortskip}{3pt}\setlength{\belowdisplayshortskip}{3pt}
\setlength{\arraycolsep}{3.5pt}\renewcommand{\arraystretch}{1.18}
\[A_{+}(1)=\begin{pmatrix}2 & 0 & 0\\0 & 2 & 0\\-2 & -2 & 2\end{pmatrix}\]
\end{minipage}%
\begin{minipage}[c]{0.495\linewidth}\centering\fontsize{10.5}{13}\selectfont
\setlength{\abovedisplayskip}{3pt}\setlength{\belowdisplayskip}{3pt}
\setlength{\abovedisplayshortskip}{3pt}\setlength{\belowdisplayshortskip}{3pt}
\setlength{\arraycolsep}{3.5pt}\renewcommand{\arraystretch}{1.18}
\[A_{-}(1)=\begin{pmatrix}2 & -1 & 0\\-1 & 2 & -1\\-1 & -2 & 2\end{pmatrix}\]
\end{minipage}%
\par\smallskip\noindent\small $h_+=12,\quad h_-=24,\quad \Omega=36.$\quad Vertices per slice: 8.
\end{minipage}\par\vfill
\noindent\begin{minipage}{\linewidth}
\noindent\textbf{Class 7}\hfill $\symbf{r}=(5,2,2)$
\par\smallskip\noindent\textcolor{rulegray}{\rule{\linewidth}{0.3pt}}
\begin{minipage}[c]{0.495\linewidth}\centering\fontsize{10.5}{13}\selectfont
\setlength{\abovedisplayskip}{3pt}\setlength{\belowdisplayskip}{3pt}
\setlength{\abovedisplayshortskip}{3pt}\setlength{\belowdisplayshortskip}{3pt}
\setlength{\arraycolsep}{3.5pt}\renewcommand{\arraystretch}{1.18}
\[N_{+}^{(7)}(z)=\begin{pmatrix}0 & 0 & 0\\0 & 0 & z\\0 & z & z\end{pmatrix}\]
\end{minipage}%
\begin{minipage}[c]{0.495\linewidth}\centering\fontsize{10.5}{13}\selectfont
\setlength{\abovedisplayskip}{3pt}\setlength{\belowdisplayskip}{3pt}
\setlength{\abovedisplayshortskip}{3pt}\setlength{\belowdisplayshortskip}{3pt}
\setlength{\arraycolsep}{3.5pt}\renewcommand{\arraystretch}{1.18}
\[N_{-}^{(7)}(z)=\begin{pmatrix}0 & z^{4} + z & z^{3} + z^{2}\\z & 0 & 0\\0 & 0 & 0\end{pmatrix}\]
\end{minipage}%
\par\noindent
\begin{minipage}[c]{0.495\linewidth}\centering\fontsize{10.5}{13}\selectfont
\setlength{\abovedisplayskip}{3pt}\setlength{\belowdisplayskip}{3pt}
\setlength{\abovedisplayshortskip}{3pt}\setlength{\belowdisplayshortskip}{3pt}
\setlength{\arraycolsep}{3.5pt}\renewcommand{\arraystretch}{1.18}
\[A_{+}(1)=\begin{pmatrix}2 & 0 & 0\\0 & 2 & -1\\0 & -1 & 1\end{pmatrix}\]
\end{minipage}%
\begin{minipage}[c]{0.495\linewidth}\centering\fontsize{10.5}{13}\selectfont
\setlength{\abovedisplayskip}{3pt}\setlength{\belowdisplayskip}{3pt}
\setlength{\abovedisplayshortskip}{3pt}\setlength{\belowdisplayshortskip}{3pt}
\setlength{\arraycolsep}{3.5pt}\renewcommand{\arraystretch}{1.18}
\[A_{-}(1)=\begin{pmatrix}2 & -2 & -2\\-1 & 2 & 0\\0 & 0 & 2\end{pmatrix}\]
\end{minipage}%
\par\smallskip\noindent\small $h_+=5,\quad h_-=7,\quad \Omega=12.$\quad Vertices per slice: 9.
\end{minipage}\par\vfill
\noindent\begin{minipage}{\linewidth}
\noindent\textbf{Class 8}\hfill $\symbf{r}=(8,2,6)$
\par\smallskip\noindent\textcolor{rulegray}{\rule{\linewidth}{0.3pt}}
\begin{minipage}[c]{0.495\linewidth}\centering\fontsize{10.5}{13}\selectfont
\setlength{\abovedisplayskip}{3pt}\setlength{\belowdisplayskip}{3pt}
\setlength{\abovedisplayshortskip}{3pt}\setlength{\belowdisplayshortskip}{3pt}
\setlength{\arraycolsep}{3.5pt}\renewcommand{\arraystretch}{1.18}
\[N_{+}^{(8)}(z)=\begin{pmatrix}0 & 0 & 0\\0 & 0 & z\\0 & z^{5} + z & 0\end{pmatrix}\]
\end{minipage}%
\begin{minipage}[c]{0.495\linewidth}\centering\fontsize{10.5}{13}\selectfont
\setlength{\abovedisplayskip}{3pt}\setlength{\belowdisplayskip}{3pt}
\setlength{\abovedisplayshortskip}{3pt}\setlength{\belowdisplayshortskip}{3pt}
\setlength{\arraycolsep}{3.5pt}\renewcommand{\arraystretch}{1.18}
\[N_{-}^{(8)}(z)=\begin{pmatrix}0 & z^{2} & z^{3} + z\\0 & 0 & 0\\z^{5} & z^{3} & 0\end{pmatrix}\]
\end{minipage}%
\par\noindent
\begin{minipage}[c]{0.495\linewidth}\centering\fontsize{10.5}{13}\selectfont
\setlength{\abovedisplayskip}{3pt}\setlength{\belowdisplayskip}{3pt}
\setlength{\abovedisplayshortskip}{3pt}\setlength{\belowdisplayshortskip}{3pt}
\setlength{\arraycolsep}{3.5pt}\renewcommand{\arraystretch}{1.18}
\[A_{+}(1)=\begin{pmatrix}2 & 0 & 0\\0 & 2 & -1\\0 & -2 & 2\end{pmatrix}\]
\end{minipage}%
\begin{minipage}[c]{0.495\linewidth}\centering\fontsize{10.5}{13}\selectfont
\setlength{\abovedisplayskip}{3pt}\setlength{\belowdisplayskip}{3pt}
\setlength{\abovedisplayshortskip}{3pt}\setlength{\belowdisplayshortskip}{3pt}
\setlength{\arraycolsep}{3.5pt}\renewcommand{\arraystretch}{1.18}
\[A_{-}(1)=\begin{pmatrix}2 & -1 & -2\\0 & 2 & 0\\-1 & -1 & 2\end{pmatrix}\]
\end{minipage}%
\par\smallskip\noindent\small $h_+=8,\quad h_-=14,\quad \Omega=22.$\quad Vertices per slice: 8.
\end{minipage}\par\vfill
\clearpage
\subsection*{Classes 9--12}
\noindent In each class, $B_\pm^{(k)}(z)=\operatorname{diag}(1+z^{r_1},1+z^{r_2},1+z^{r_3})-N_\pm^{(k)}(z)$.\par\medskip
\noindent\begin{minipage}{\linewidth}
\noindent\textbf{Class 9}\hfill $\symbf{r}=(5,2,3)$
\par\smallskip\noindent\textcolor{rulegray}{\rule{\linewidth}{0.3pt}}
\begin{minipage}[c]{0.495\linewidth}\centering\fontsize{10.5}{13}\selectfont
\setlength{\abovedisplayskip}{3pt}\setlength{\belowdisplayskip}{3pt}
\setlength{\abovedisplayshortskip}{3pt}\setlength{\belowdisplayshortskip}{3pt}
\setlength{\arraycolsep}{3.5pt}\renewcommand{\arraystretch}{1.18}
\[N_{+}^{(9)}(z)=\begin{pmatrix}0 & 0 & 0\\0 & 0 & z\\0 & z^{2} + z & 0\end{pmatrix}\]
\end{minipage}%
\begin{minipage}[c]{0.495\linewidth}\centering\fontsize{10.5}{13}\selectfont
\setlength{\abovedisplayskip}{3pt}\setlength{\belowdisplayskip}{3pt}
\setlength{\abovedisplayshortskip}{3pt}\setlength{\belowdisplayshortskip}{3pt}
\setlength{\arraycolsep}{3.5pt}\renewcommand{\arraystretch}{1.18}
\[N_{-}^{(9)}(z)=\begin{pmatrix}0 & z^{2} & z^{3} + z\\0 & z & 0\\z^{2} & 0 & 0\end{pmatrix}\]
\end{minipage}%
\par\noindent
\begin{minipage}[c]{0.495\linewidth}\centering\fontsize{10.5}{13}\selectfont
\setlength{\abovedisplayskip}{3pt}\setlength{\belowdisplayskip}{3pt}
\setlength{\abovedisplayshortskip}{3pt}\setlength{\belowdisplayshortskip}{3pt}
\setlength{\arraycolsep}{3.5pt}\renewcommand{\arraystretch}{1.18}
\[A_{+}(1)=\begin{pmatrix}2 & 0 & 0\\0 & 2 & -1\\0 & -2 & 2\end{pmatrix}\]
\end{minipage}%
\begin{minipage}[c]{0.495\linewidth}\centering\fontsize{10.5}{13}\selectfont
\setlength{\abovedisplayskip}{3pt}\setlength{\belowdisplayskip}{3pt}
\setlength{\abovedisplayshortskip}{3pt}\setlength{\belowdisplayshortskip}{3pt}
\setlength{\arraycolsep}{3.5pt}\renewcommand{\arraystretch}{1.18}
\[A_{-}(1)=\begin{pmatrix}2 & -1 & -2\\0 & 1 & 0\\-1 & 0 & 2\end{pmatrix}\]
\end{minipage}%
\par\smallskip\noindent\small $h_+=5,\quad h_-=8,\quad \Omega=13.$\quad Vertices per slice: 10.
\end{minipage}\par\vfill
\noindent\begin{minipage}{\linewidth}
\noindent\textbf{Class 10}\hfill $\symbf{r}=(8,2,6)$
\par\smallskip\noindent\textcolor{rulegray}{\rule{\linewidth}{0.3pt}}
\begin{minipage}[c]{0.495\linewidth}\centering\fontsize{10.5}{13}\selectfont
\setlength{\abovedisplayskip}{3pt}\setlength{\belowdisplayskip}{3pt}
\setlength{\abovedisplayshortskip}{3pt}\setlength{\belowdisplayshortskip}{3pt}
\setlength{\arraycolsep}{3.5pt}\renewcommand{\arraystretch}{1.18}
\[N_{+}^{(10)}(z)=\begin{pmatrix}0 & 0 & 0\\0 & 0 & z\\0 & z^{5} + z & 0\end{pmatrix}\]
\end{minipage}%
\begin{minipage}[c]{0.495\linewidth}\centering\fontsize{10.5}{13}\selectfont
\setlength{\abovedisplayskip}{3pt}\setlength{\belowdisplayskip}{3pt}
\setlength{\abovedisplayshortskip}{3pt}\setlength{\belowdisplayshortskip}{3pt}
\setlength{\arraycolsep}{3.5pt}\renewcommand{\arraystretch}{1.18}
\[N_{-}^{(10)}(z)=\begin{pmatrix}0 & z^{7} + z & z^{6} + z^{2}\\z & 0 & 0\\0 & z^{3} & 0\end{pmatrix}\]
\end{minipage}%
\par\noindent
\begin{minipage}[c]{0.495\linewidth}\centering\fontsize{10.5}{13}\selectfont
\setlength{\abovedisplayskip}{3pt}\setlength{\belowdisplayskip}{3pt}
\setlength{\abovedisplayshortskip}{3pt}\setlength{\belowdisplayshortskip}{3pt}
\setlength{\arraycolsep}{3.5pt}\renewcommand{\arraystretch}{1.18}
\[A_{+}(1)=\begin{pmatrix}2 & 0 & 0\\0 & 2 & -1\\0 & -2 & 2\end{pmatrix}\]
\end{minipage}%
\begin{minipage}[c]{0.495\linewidth}\centering\fontsize{10.5}{13}\selectfont
\setlength{\abovedisplayskip}{3pt}\setlength{\belowdisplayskip}{3pt}
\setlength{\abovedisplayshortskip}{3pt}\setlength{\belowdisplayshortskip}{3pt}
\setlength{\arraycolsep}{3.5pt}\renewcommand{\arraystretch}{1.18}
\[A_{-}(1)=\begin{pmatrix}2 & -2 & -2\\-1 & 2 & 0\\0 & -1 & 2\end{pmatrix}\]
\end{minipage}%
\par\smallskip\noindent\small $h_+=8,\quad h_-=16,\quad \Omega=24.$\quad Vertices per slice: 8.
\end{minipage}\par\vfill
\noindent\begin{minipage}{\linewidth}
\noindent\textbf{Class 11}\hfill $\symbf{r}=(2,2,8)$
\par\smallskip\noindent\textcolor{rulegray}{\rule{\linewidth}{0.3pt}}
\begin{minipage}[c]{0.495\linewidth}\centering\fontsize{10.5}{13}\selectfont
\setlength{\abovedisplayskip}{3pt}\setlength{\belowdisplayskip}{3pt}
\setlength{\abovedisplayshortskip}{3pt}\setlength{\belowdisplayshortskip}{3pt}
\setlength{\arraycolsep}{3.5pt}\renewcommand{\arraystretch}{1.18}
\[N_{+}^{(11)}(z)=\begin{pmatrix}0 & 0 & 0\\0 & 0 & z\\z^{6} + z^{2} & z^{7} + z & 0\end{pmatrix}\]
\end{minipage}%
\begin{minipage}[c]{0.495\linewidth}\centering\fontsize{10.5}{13}\selectfont
\setlength{\abovedisplayskip}{3pt}\setlength{\belowdisplayskip}{3pt}
\setlength{\abovedisplayshortskip}{3pt}\setlength{\belowdisplayshortskip}{3pt}
\setlength{\arraycolsep}{3.5pt}\renewcommand{\arraystretch}{1.18}
\[N_{-}^{(11)}(z)=\begin{pmatrix}0 & z & 0\\z & 0 & 0\\z^{4} & 0 & 0\end{pmatrix}\]
\end{minipage}%
\par\noindent
\begin{minipage}[c]{0.495\linewidth}\centering\fontsize{10.5}{13}\selectfont
\setlength{\abovedisplayskip}{3pt}\setlength{\belowdisplayskip}{3pt}
\setlength{\abovedisplayshortskip}{3pt}\setlength{\belowdisplayshortskip}{3pt}
\setlength{\arraycolsep}{3.5pt}\renewcommand{\arraystretch}{1.18}
\[A_{+}(1)=\begin{pmatrix}2 & 0 & 0\\0 & 2 & -1\\-2 & -2 & 2\end{pmatrix}\]
\end{minipage}%
\begin{minipage}[c]{0.495\linewidth}\centering\fontsize{10.5}{13}\selectfont
\setlength{\abovedisplayskip}{3pt}\setlength{\belowdisplayskip}{3pt}
\setlength{\abovedisplayshortskip}{3pt}\setlength{\belowdisplayshortskip}{3pt}
\setlength{\arraycolsep}{3.5pt}\renewcommand{\arraystretch}{1.18}
\[A_{-}(1)=\begin{pmatrix}2 & -1 & 0\\-1 & 2 & 0\\-1 & 0 & 2\end{pmatrix}\]
\end{minipage}%
\par\smallskip\noindent\small $h_+=10,\quad h_-=8,\quad \Omega=18.$\quad Vertices per slice: 6.
\end{minipage}\par\vfill
\noindent\begin{minipage}{\linewidth}
\noindent\textbf{Class 12}\hfill $\symbf{r}=(2,2,12)$
\par\smallskip\noindent\textcolor{rulegray}{\rule{\linewidth}{0.3pt}}
\begin{minipage}[c]{0.495\linewidth}\centering\fontsize{10.5}{13}\selectfont
\setlength{\abovedisplayskip}{3pt}\setlength{\belowdisplayskip}{3pt}
\setlength{\abovedisplayshortskip}{3pt}\setlength{\belowdisplayshortskip}{3pt}
\setlength{\arraycolsep}{3.5pt}\renewcommand{\arraystretch}{1.18}
\[N_{+}^{(12)}(z)=\begin{pmatrix}0 & 0 & 0\\0 & 0 & z\\z^{10} + z^{6} + z^{2} & z^{11} + z & 0\end{pmatrix}\]
\end{minipage}%
\begin{minipage}[c]{0.495\linewidth}\centering\fontsize{10.5}{13}\selectfont
\setlength{\abovedisplayskip}{3pt}\setlength{\belowdisplayskip}{3pt}
\setlength{\abovedisplayshortskip}{3pt}\setlength{\belowdisplayshortskip}{3pt}
\setlength{\arraycolsep}{3.5pt}\renewcommand{\arraystretch}{1.18}
\[N_{-}^{(12)}(z)=\begin{pmatrix}0 & z & 0\\z & 0 & 0\\z^{8} + z^{4} & 0 & z^{6}\end{pmatrix}\]
\end{minipage}%
\par\noindent
\begin{minipage}[c]{0.495\linewidth}\centering\fontsize{10.5}{13}\selectfont
\setlength{\abovedisplayskip}{3pt}\setlength{\belowdisplayskip}{3pt}
\setlength{\abovedisplayshortskip}{3pt}\setlength{\belowdisplayshortskip}{3pt}
\setlength{\arraycolsep}{3.5pt}\renewcommand{\arraystretch}{1.18}
\[A_{+}(1)=\begin{pmatrix}2 & 0 & 0\\0 & 2 & -1\\-3 & -2 & 2\end{pmatrix}\]
\end{minipage}%
\begin{minipage}[c]{0.495\linewidth}\centering\fontsize{10.5}{13}\selectfont
\setlength{\abovedisplayskip}{3pt}\setlength{\belowdisplayskip}{3pt}
\setlength{\abovedisplayshortskip}{3pt}\setlength{\belowdisplayshortskip}{3pt}
\setlength{\arraycolsep}{3.5pt}\renewcommand{\arraystretch}{1.18}
\[A_{-}(1)=\begin{pmatrix}2 & -1 & 0\\-1 & 2 & 0\\-2 & 0 & 1\end{pmatrix}\]
\end{minipage}%
\par\smallskip\noindent\small $h_+=14,\quad h_-=18,\quad \Omega=32.$\quad Vertices per slice: 8.
\end{minipage}\par\vfill
\clearpage
\subsection*{Classes 13--16}
\noindent In each class, $B_\pm^{(k)}(z)=\operatorname{diag}(1+z^{r_1},1+z^{r_2},1+z^{r_3})-N_\pm^{(k)}(z)$.\par\medskip
\noindent\begin{minipage}{\linewidth}
\noindent\textbf{Class 13}\hfill $\symbf{r}=(3,2,2)$
\par\smallskip\noindent\textcolor{rulegray}{\rule{\linewidth}{0.3pt}}
\begin{minipage}[c]{0.495\linewidth}\centering\fontsize{10.5}{13}\selectfont
\setlength{\abovedisplayskip}{3pt}\setlength{\belowdisplayskip}{3pt}
\setlength{\abovedisplayshortskip}{3pt}\setlength{\belowdisplayshortskip}{3pt}
\setlength{\arraycolsep}{3.5pt}\renewcommand{\arraystretch}{1.18}
\[N_{+}^{(13)}(z)=\begin{pmatrix}0 & 0 & 0\\0 & z & 0\\0 & 0 & z\end{pmatrix}\]
\end{minipage}%
\begin{minipage}[c]{0.495\linewidth}\centering\fontsize{10.5}{13}\selectfont
\setlength{\abovedisplayskip}{3pt}\setlength{\belowdisplayskip}{3pt}
\setlength{\abovedisplayshortskip}{3pt}\setlength{\belowdisplayshortskip}{3pt}
\setlength{\arraycolsep}{3.5pt}\renewcommand{\arraystretch}{1.18}
\[N_{-}^{(13)}(z)=\begin{pmatrix}0 & 0 & z^{2} + z\\0 & 0 & z\\z & z & 0\end{pmatrix}\]
\end{minipage}%
\par\noindent
\begin{minipage}[c]{0.495\linewidth}\centering\fontsize{10.5}{13}\selectfont
\setlength{\abovedisplayskip}{3pt}\setlength{\belowdisplayskip}{3pt}
\setlength{\abovedisplayshortskip}{3pt}\setlength{\belowdisplayshortskip}{3pt}
\setlength{\arraycolsep}{3.5pt}\renewcommand{\arraystretch}{1.18}
\[A_{+}(1)=\begin{pmatrix}2 & 0 & 0\\0 & 1 & 0\\0 & 0 & 1\end{pmatrix}\]
\end{minipage}%
\begin{minipage}[c]{0.495\linewidth}\centering\fontsize{10.5}{13}\selectfont
\setlength{\abovedisplayskip}{3pt}\setlength{\belowdisplayskip}{3pt}
\setlength{\abovedisplayshortskip}{3pt}\setlength{\belowdisplayshortskip}{3pt}
\setlength{\arraycolsep}{3.5pt}\renewcommand{\arraystretch}{1.18}
\[A_{-}(1)=\begin{pmatrix}2 & 0 & -2\\0 & 2 & -1\\-1 & -1 & 2\end{pmatrix}\]
\end{minipage}%
\par\smallskip\noindent\small $h_+=3,\quad h_-=7,\quad \Omega=10.$\quad Vertices per slice: 7.
\end{minipage}\par\vfill
\noindent\begin{minipage}{\linewidth}
\noindent\textbf{Class 14}\hfill $\symbf{r}=(2,12,10)$
\par\smallskip\noindent\textcolor{rulegray}{\rule{\linewidth}{0.3pt}}
\begin{minipage}[c]{0.495\linewidth}\centering\fontsize{10.5}{13}\selectfont
\setlength{\abovedisplayskip}{3pt}\setlength{\belowdisplayskip}{3pt}
\setlength{\abovedisplayshortskip}{3pt}\setlength{\belowdisplayshortskip}{3pt}
\setlength{\arraycolsep}{3.5pt}\renewcommand{\arraystretch}{1.18}
\[N_{+}^{(14)}(z)=\begin{pmatrix}0 & 0 & 0\\z^{2} & 0 & z^{3} + z\\z^{7} + z^{3} & z^{9} & 0\end{pmatrix}\]
\end{minipage}%
\begin{minipage}[c]{0.495\linewidth}\centering\fontsize{10.5}{13}\selectfont
\setlength{\abovedisplayskip}{3pt}\setlength{\belowdisplayskip}{3pt}
\setlength{\abovedisplayshortskip}{3pt}\setlength{\belowdisplayshortskip}{3pt}
\setlength{\arraycolsep}{3.5pt}\renewcommand{\arraystretch}{1.18}
\[N_{-}^{(14)}(z)=\begin{pmatrix}0 & 0 & z\\0 & z^{6} & 0\\z^{9} + z^{5} + z & 0 & 0\end{pmatrix}\]
\end{minipage}%
\par\noindent
\begin{minipage}[c]{0.495\linewidth}\centering\fontsize{10.5}{13}\selectfont
\setlength{\abovedisplayskip}{3pt}\setlength{\belowdisplayskip}{3pt}
\setlength{\abovedisplayshortskip}{3pt}\setlength{\belowdisplayshortskip}{3pt}
\setlength{\arraycolsep}{3.5pt}\renewcommand{\arraystretch}{1.18}
\[A_{+}(1)=\begin{pmatrix}2 & 0 & 0\\-1 & 2 & -2\\-2 & -1 & 2\end{pmatrix}\]
\end{minipage}%
\begin{minipage}[c]{0.495\linewidth}\centering\fontsize{10.5}{13}\selectfont
\setlength{\abovedisplayskip}{3pt}\setlength{\belowdisplayskip}{3pt}
\setlength{\abovedisplayshortskip}{3pt}\setlength{\belowdisplayshortskip}{3pt}
\setlength{\arraycolsep}{3.5pt}\renewcommand{\arraystretch}{1.18}
\[A_{-}(1)=\begin{pmatrix}2 & 0 & -1\\0 & 1 & 0\\-3 & 0 & 2\end{pmatrix}\]
\end{minipage}%
\par\smallskip\noindent\small $h_+=22,\quad h_-=18,\quad \Omega=40.$\quad Vertices per slice: 12.
\end{minipage}\par\vfill
\noindent\begin{minipage}{\linewidth}
\noindent\textbf{Class 15}\hfill $\symbf{r}=(2,2,2)$
\par\smallskip\noindent\textcolor{rulegray}{\rule{\linewidth}{0.3pt}}
\begin{minipage}[c]{0.495\linewidth}\centering\fontsize{10.5}{13}\selectfont
\setlength{\abovedisplayskip}{3pt}\setlength{\belowdisplayskip}{3pt}
\setlength{\abovedisplayshortskip}{3pt}\setlength{\belowdisplayshortskip}{3pt}
\setlength{\arraycolsep}{3.5pt}\renewcommand{\arraystretch}{1.18}
\[N_{+}^{(15)}(z)=\begin{pmatrix}0 & 0 & z\\0 & 0 & z\\z & z & 0\end{pmatrix}\]
\end{minipage}%
\begin{minipage}[c]{0.495\linewidth}\centering\fontsize{10.5}{13}\selectfont
\setlength{\abovedisplayskip}{3pt}\setlength{\belowdisplayskip}{3pt}
\setlength{\abovedisplayshortskip}{3pt}\setlength{\belowdisplayshortskip}{3pt}
\setlength{\arraycolsep}{3.5pt}\renewcommand{\arraystretch}{1.18}
\[N_{-}^{(15)}(z)=\begin{pmatrix}0 & z & 0\\z & 0 & 0\\0 & 0 & z\end{pmatrix}\]
\end{minipage}%
\par\noindent
\begin{minipage}[c]{0.495\linewidth}\centering\fontsize{10.5}{13}\selectfont
\setlength{\abovedisplayskip}{3pt}\setlength{\belowdisplayskip}{3pt}
\setlength{\abovedisplayshortskip}{3pt}\setlength{\belowdisplayshortskip}{3pt}
\setlength{\arraycolsep}{3.5pt}\renewcommand{\arraystretch}{1.18}
\[A_{+}(1)=\begin{pmatrix}2 & 0 & -1\\0 & 2 & -1\\-1 & -1 & 2\end{pmatrix}\]
\end{minipage}%
\begin{minipage}[c]{0.495\linewidth}\centering\fontsize{10.5}{13}\selectfont
\setlength{\abovedisplayskip}{3pt}\setlength{\belowdisplayskip}{3pt}
\setlength{\abovedisplayshortskip}{3pt}\setlength{\belowdisplayshortskip}{3pt}
\setlength{\arraycolsep}{3.5pt}\renewcommand{\arraystretch}{1.18}
\[A_{-}(1)=\begin{pmatrix}2 & -1 & 0\\-1 & 2 & 0\\0 & 0 & 1\end{pmatrix}\]
\end{minipage}%
\par\smallskip\noindent\small $h_+=4,\quad h_-=3,\quad \Omega=7.$\quad Vertices per slice: 6.
\end{minipage}\par\vfill
\noindent\begin{minipage}{\linewidth}
\noindent\textbf{Class 16}\hfill $\symbf{r}=(2,2,2)$
\par\smallskip\noindent\textcolor{rulegray}{\rule{\linewidth}{0.3pt}}
\begin{minipage}[c]{0.495\linewidth}\centering\fontsize{10.5}{13}\selectfont
\setlength{\abovedisplayskip}{3pt}\setlength{\belowdisplayskip}{3pt}
\setlength{\abovedisplayshortskip}{3pt}\setlength{\belowdisplayshortskip}{3pt}
\setlength{\arraycolsep}{3.5pt}\renewcommand{\arraystretch}{1.18}
\[N_{+}^{(16)}(z)=\begin{pmatrix}0 & 0 & z\\0 & 0 & z\\z & z & 0\end{pmatrix}\]
\end{minipage}%
\begin{minipage}[c]{0.495\linewidth}\centering\fontsize{10.5}{13}\selectfont
\setlength{\abovedisplayskip}{3pt}\setlength{\belowdisplayskip}{3pt}
\setlength{\abovedisplayshortskip}{3pt}\setlength{\belowdisplayshortskip}{3pt}
\setlength{\arraycolsep}{3.5pt}\renewcommand{\arraystretch}{1.18}
\[N_{-}^{(16)}(z)=\begin{pmatrix}z & 0 & 0\\0 & z & 0\\0 & 0 & z\end{pmatrix}\]
\end{minipage}%
\par\noindent
\begin{minipage}[c]{0.495\linewidth}\centering\fontsize{10.5}{13}\selectfont
\setlength{\abovedisplayskip}{3pt}\setlength{\belowdisplayskip}{3pt}
\setlength{\abovedisplayshortskip}{3pt}\setlength{\belowdisplayshortskip}{3pt}
\setlength{\arraycolsep}{3.5pt}\renewcommand{\arraystretch}{1.18}
\[A_{+}(1)=\begin{pmatrix}2 & 0 & -1\\0 & 2 & -1\\-1 & -1 & 2\end{pmatrix}\]
\end{minipage}%
\begin{minipage}[c]{0.495\linewidth}\centering\fontsize{10.5}{13}\selectfont
\setlength{\abovedisplayskip}{3pt}\setlength{\belowdisplayskip}{3pt}
\setlength{\abovedisplayshortskip}{3pt}\setlength{\belowdisplayshortskip}{3pt}
\setlength{\arraycolsep}{3.5pt}\renewcommand{\arraystretch}{1.18}
\[A_{-}(1)=\begin{pmatrix}1 & 0 & 0\\0 & 1 & 0\\0 & 0 & 1\end{pmatrix}\]
\end{minipage}%
\par\smallskip\noindent\small $h_+=4,\quad h_-=3,\quad \Omega=14.$\quad Vertices per slice: 6.
\end{minipage}\par\vfill


\clearpage
\appendix
\section{Reproduction and certificate files}\label{app:reproduce}

The classification programs and certificates are under
\code{research/rank3/}. The checked runtime is Python 3.10.6,
SymPy 1.14.0, NumPy 2.2.6 and Z3 5.1.0. Z3 is installed locally under
the ignored \code{_deps/} directory.

Run the following commands from the repository root:
\begin{quote}\small\ttfamily
python -m pip install --target research/rank3/\_deps\\
\hspace*{1em}z3-solver==5.1.0.0\quad\normalfont(on the same command line)\ttfamily\\[3pt]
python research/rank3/classify\_rank3.py\\
python research/rank3/enumerate\_constants.py\\
python research/rank3/solve\_lifts.py --timeout 2000\\
\hspace*{1em}--optimize --dynamics\quad\normalfont(on the same command line)\ttfamily\\[3pt]
python research/rank3/classify\_lift\_spaces.py\\
python research/rank3/verify\_classification.py\\
python research/rank3/slice\_equivalence.py
\end{quote}

A timeout is reported as UNKNOWN and cannot pass the verifier as an
exclusion. The successful run has no UNKNOWN results. Table~\ref{tab:files}
describes the certificate files.

\begin{table}[ht]
\centering\small\renewcommand{\arraystretch}{1.2}
\begin{tabular}{@{}p{0.45\linewidth}p{0.50\linewidth}@{}}
\toprule
File & Content\\\midrule
\code{constant_candidates.json} & All 180 constant pairs and positive vectors.\\
\code{lift_feasibility.jsonl} & Solver outcomes, sixteen witnesses and periodicity certificates.\\
\code{lift_spaces.json} & Exact row-reduced equations of all sixteen families.\\
\code{verification.json} & The replayed UNSAT answers, query hashes and runtime versions.\\
\code{smt_queries/} & All 180 unbounded SMT-LIB2 formulas.\\
\code{slice_signatures.json} & Complete encodings distinguishing the slice classes.\\
\bottomrule
\end{tabular}
\caption{Certificate files under \code{research/rank3/}.}\label{tab:files}
\end{table}

The earlier paper \cite[Table 3]{MizunoTData} gives seven rank-three
examples. In their original order, they correspond to classes
$1,16,13,10,2,4,7$ here. The other nine classes are additional to that
table.

\paragraph{Building this document.}
The overview and matrix catalogue are generated directly from the
verified JSON by \code{make_pdf_tables.py}. The source is
\code{rank3-classification.tex}; the PDF is built with LuaLaTeX and
\code{latexmk}. All matrix coefficients in the catalogue are generated
from the exact witnesses.

\begin{thebibliography}{9}
\bibitem{MizunoRankTwo}
Y.~Mizuno, \emph{Periodic Y-systems and Nahm sums: the rank 2 case},
SIGMA \textbf{21} (2025), 094, 17 pp.
\href{https://doi.org/10.3842/SIGMA.2025.094}{doi:10.3842/SIGMA.2025.094}.
\href{https://arxiv.org/abs/2301.13239v2}{arXiv:2301.13239v2}.

\bibitem{MizunoTData}
Y.~Mizuno, \emph{Difference equations arising from cluster algebras},
\href{https://arxiv.org/abs/1912.05710v2}{arXiv:1912.05710v2} (2020).
\end{thebibliography}
\end{document}
